\chapter{Propositional Logic: Syntax}\label{ch:propsyntax}

Syntax is about which strings of symbols are legal formulas, with no reference yet to
truth or meaning.

\section{The alphabet}

The language of propositional logic, $\mathcal{L}^p$, uses the following symbols:
\begin{itemize}
  \item \emph{propositional variables} (atoms): $p, q, r, \dots$, possibly with
        subscripts $p_1, p_2, \dots$; the set of atoms is written $\Prop$ (countable);
  \item \emph{connectives}: $\neg$ (negation), $\land$ (conjunction), $\lor$ (disjunction),
        $\to$ (implication), $\leftrightarrow$ (biconditional);
  \item \emph{punctuation}: the parentheses $($ and $)$.
\end{itemize}
Some presentations also include the constants $\top$ and $\bot$ as atomic formulas.

An \emph{expression} is any finite string of these symbols. Most expressions, such as
$)p\neg\land($, are garbage. We now pick out the meaningful ones.

\section{Well-formed formulas}

\begin{definition}[Well-formed formula]
The set $\Fm(\mathcal{L}^p)$ of \emph{well-formed formulas} (wffs, or just formulas) is
the smallest set of expressions such that
\begin{enumerate}
  \item every atom $p \in \Prop$ is a formula (an \emph{atomic formula});
  \item if $A$ is a formula, then $(\neg A)$ is a formula;
  \item if $A$ and $B$ are formulas and $\star \in \{\land,\lor,\to,\leftrightarrow\}$,
        then $(A \star B)$ is a formula.
\end{enumerate}
\end{definition}

``Smallest'' means nothing is a formula unless it is forced to be by finitely many
applications of the rules. Equivalently: an expression is a formula iff it has a
\emph{formation sequence}, a finite list of expressions each of which is an atom or is
obtained from earlier ones by rule 2 or 3, ending with the expression.

\begin{example}
$((p \land q) \to (\neg r))$ is a formula, with formation sequence
\[
  p,\quad q,\quad r,\quad (p\land q),\quad (\neg r),\quad ((p\land q)\to(\neg r)).
\]
The expressions $p\land q$ (missing outer parentheses), $(p \neg q)$ and $(p \land)$ are
not formulas under the strict definition.
\end{example}

\subsection{Conventions for dropping parentheses}

Fully parenthesised formulas are hard to read, so we adopt conventions (and always
mean the fully parenthesised formula):
\begin{itemize}
  \item omit the outermost pair of parentheses;
  \item precedence, from tightest to loosest: $\neg$, $\land$, $\lor$, $\to$,
        $\leftrightarrow$;
  \item $\to$ associates to the \emph{right}: $p \to q \to r$ means $p\to(q\to r)$;
  \item $\land$ and $\lor$ are usually written associating left; since they are
        associative semantically the choice rarely matters.
\end{itemize}
So $\neg p \land q \to r \lor s$ abbreviates $(((\neg p) \land q) \to (r \lor s))$.

\section{Structural induction}

Because formulas are defined inductively, we prove properties of \emph{all} formulas by
induction that follows the definition.

\begin{theorem}[Principle of structural induction]
Let $R$ be a property of expressions. Suppose
\begin{enumerate}
  \item (base case) every atom has property $R$;
  \item (inductive step for $\neg$) if $A$ has $R$, then $(\neg A)$ has $R$;
  \item (inductive step for binary connectives) if $A$ and $B$ have $R$, then so does
        $(A \star B)$ for each $\star\in\{\land,\lor,\to,\leftrightarrow\}$.
\end{enumerate}
Then every well-formed formula has property $R$.
\end{theorem}

\begin{proof}[Justification]
Let $S$ be the set of formulas with property $R$. The hypotheses say $S$ contains all
atoms and is closed under the formation rules. Since $\Fm$ is the \emph{smallest} such
set, $\Fm \subseteq S$. (Alternatively, do ordinary strong induction on the length of a
formation sequence, or on the number of connectives.)
\end{proof}

\begin{example}[Balanced parentheses]
\emph{Claim}: every formula has the same number of left and right parentheses.

\emph{Proof.} Let $l(A)$, $r(A)$ be the numbers of left and right parentheses in $A$.
\begin{itemize}
  \item Base: an atom $p$ has $l(p)=r(p)=0$.
  \item $(\neg A)$: $l = l(A)+1 = r(A)+1 = r$ by the induction hypothesis.
  \item $(A\star B)$: $l = l(A)+l(B)+1 = r(A)+r(B)+1 = r$.
\end{itemize}
By structural induction the claim holds for every formula. \qed
\end{example}

\begin{example}[Proper prefixes]
\emph{Claim}: every non-empty proper prefix of a formula has strictly more left than
right parentheses (in particular, no proper prefix of a formula is a formula).

\emph{Sketch.} Atoms have no non-empty proper prefixes. For $(\neg A)$, a proper prefix is
$($, or $(\neg$, or $(\neg A'$ where $A'$ is a prefix of $A$ (a proper prefix, or $A$
itself); in each case the extra leading $($ is unmatched, using the IH for proper
prefixes of $A$ and the balance lemma for $A$ itself. The case $(A\star B)$ is similar.
\end{example}

\begin{example}[Counting]
Let $c(A)$ be the number of occurrences of binary connectives and $a(A)$ the number of
occurrences of atoms. If $A$ contains no $\neg$, then $a(A) = c(A)+1$. Base: $1 = 0+1$.
Step: $a(A\star B) = a(A)+a(B) = c(A)+1+c(B)+1 = c(A\star B)+1$.
\end{example}

\section{Unique readability}

Could a formula be built in two different ways, and so have two different meanings? For
example, is $(p\land q\to r)$ ambiguous? Under the official, fully parenthesised
definition, no.

\begin{theorem}[Unique readability]
Every formula is of exactly one of the following forms:
\begin{enumerate}
  \item an atom;
  \item $(\neg A)$ for a unique formula $A$;
  \item $(A\star B)$ for a unique connective $\star$ and unique formulas $A$, $B$.
\end{enumerate}
\end{theorem}

\begin{proof}[Proof sketch]
The first symbol of a formula is an atom (case 1) or $($. If it starts with $(\neg$ then
it must come from rule 2, and $A$ is determined by removing the first two and last
symbols. Otherwise suppose $(A\star B) = (A'\star' B')$. Then $A$ and $A'$ both start at
position 2. If $A \neq A'$, one is a proper prefix of the other, so a proper prefix of a
formula is a formula, contradicting the prefix lemma. Hence $A=A'$, so $\star=\star'$
and $B=B'$.
\end{proof}

Unique readability is what lets us define functions on formulas by \emph{structural
recursion} (e.g.\ the truth value of a formula, its height, its set of subformulas) and
know they are well defined.

\begin{definition}[Main connective, subformula]
In case 2 or 3 above, $\neg$ (respectively $\star$) is the \emph{main connective} of the
formula. The set of \emph{subformulas} is defined recursively:
$\mathrm{sub}(p)=\{p\}$, $\mathrm{sub}((\neg A))=\{(\neg A)\}\cup\mathrm{sub}(A)$, and
$\mathrm{sub}((A\star B)) = \{(A\star B)\}\cup\mathrm{sub}(A)\cup\mathrm{sub}(B)$.
\end{definition}

\section{Parse trees}

A formula can be drawn as a \emph{parse tree}: leaves are atoms, each internal node is
labelled with a connective, $\neg$ nodes have one child and binary nodes have two. The
root is the main connective. Unique readability says each formula has exactly one parse
tree.

\begin{example}
The parse tree of $\neg p \land q \to (r \lor \neg p)$, i.e.\
$(((\neg p)\land q)\to(r\lor(\neg p)))$:
\begin{center}
\begin{tikzpicture}[level distance=11mm,
  level 1/.style={sibling distance=36mm},
  level 2/.style={sibling distance=16mm},
  level 3/.style={sibling distance=10mm}]
\node {$\to$}
  child {node {$\land$}
    child {node {$\neg$} child {node {$p$}}}
    child {node {$q$}}}
  child {node {$\lor$}
    child {node {$r$}}
    child {node {$\neg$} child {node {$p$}}}};
\end{tikzpicture}
\end{center}
Every subtree corresponds to a subformula. The \emph{height} of the formula is the height
of the tree (here $3$, counting edges on the longest root-to-leaf path).
\end{example}

\begin{aside}
Checking whether a string is a well-formed formula is a parsing problem, like in a
compiler. A simple algorithm: if the string is an atom, accept; if it starts with
$(\neg$ and ends with $)$, recurse on the middle; otherwise, scan from the second symbol
keeping a parenthesis counter and find the first point where the counter returns to $0$
after a complete formula; the next symbol must be a binary connective, and recurse on
both sides. The prefix lemma guarantees this split point is the only possible one.
\end{aside}

\section{Recursive definitions on formulas}

Some standard functions defined by structural recursion:
\begin{align*}
  \mathrm{height}(p) &= 0, &
  \mathrm{height}(\neg A) &= 1 + \mathrm{height}(A), \\
  && \mathrm{height}(A\star B) &= 1+\max(\mathrm{height}(A),\mathrm{height}(B)).\\
  \mathrm{atoms}(p) &= \{p\}, &
  \mathrm{atoms}(\neg A) &= \mathrm{atoms}(A), \\
  && \mathrm{atoms}(A\star B) &= \mathrm{atoms}(A)\cup\mathrm{atoms}(B).
\end{align*}
Proofs about such functions are then by structural induction. For instance, one can show
a formula of height $h$ has at most $2^{h+1}-1$ symbols excluding parentheses.

\section*{Summary}
\begin{itemize}
  \item Formulas are the smallest set containing the atoms and closed under
        $(\neg A)$ and $(A\star B)$.
  \item Structural induction: prove for atoms, then show each formation rule preserves
        the property.
  \item Lemmas: parentheses balance; proper prefixes have excess left parentheses.
  \item Unique readability: every formula decomposes in exactly one way, so it has one
        parse tree and a unique main connective, and recursive definitions are well
        defined.
  \item Precedence conventions ($\neg,\land,\lor,\to,\leftrightarrow$; $\to$ is
        right-associative) are only abbreviations.
\end{itemize}
